\documentclass[10pt,a4paper]{amsart}
\usepackage[margin=19mm]{geometry}
\usepackage{amsmath,amssymb,mathtools}
\usepackage[scr=rsfs,cal=euler]{mathalfa}
\usepackage{xfrac}
\usepackage[unicode,hidelinks,bookmarks=false]{hyperref}
\newcommand{\ie}{i.e.\ }
\newcommand{\cf}{cf.\ }
\newcommand{\resp}{respectively\ }
\newcommand{\Subsection}{\subsection}
\providecommand{\nobreakdash}{\nobreak\hbox{-}}
\setlength{\emergencystretch}{2em}
\multlinegap0pt
\usepackage{bm}
\usepackage{bbm}
\usepackage{dsfont}
\usepackage{xspace}
\usepackage{cancel}
\usepackage{mathtools}
\usepackage{amsthm}
\makeatletter
\@ifundefined{theorem}{\newtheorem{theorem}{Theorem}}{}
\@ifundefined{proposition}{\newtheorem{proposition}[theorem]{Proposition}}{}
\makeatother
\makeatletter
\renewcommand{\P}{\mathbb{P}} % Polynomials
\newcommand{\de}{\mathrm{d}}
\expandafter\ifx\csname proof\endcsname\relax
\newenvironment{proof}[1][Proof]{\noindent\textbf{#1.} }{\ \rule{0.5em}{0.5em}}
\fi
\makeatother
\title[Bivariate polynomial histopolation techniques on Padua, Feke]{Bivariate polynomial histopolation techniques on Padua, Fekete and Leja triangles}
\author{Ludovico Bruni Bruno, Francesco Dell'Accio, Wolfgang Erb, Federico Nudo}
\date{Source: arXiv:2506.03025v1 (CC BY 4.0)}
\begin{document}
\maketitle

\noindent\emph{Source and licence.} Bivariate polynomial histopolation techniques on Padua, Fekete and Leja triangles. Version arXiv:2506.03025v1 (CC BY 4.0), distributed under the Creative Commons Attribution 4.0 International licence, as recorded in the arXiv licence metadata for this exact version and shown on the abstract page https://arxiv.org/abs/2506.03025v1. Full rights notice: licenses/arxiv-2506.03025-CC-BY-4.0.txt.\par\smallskip

\noindent\textit{Editorial scope.} Counted unit: the Proposition whose source label is \texttt{Thm:EstimateOfNorm} in the article (source0.tex lines 570--580, source label \texttt{Thm:EstimateOfNorm}), delivered together with the article's own complete original proof of it (source0.tex lines 583--646, one \texttt{proof} environment of 64 lines). No mathematics is written, repaired or re-derived: statement and proof are the article's own text line for line. Required context retained uncounted: a1eq (lines 532--534); a2eq (lines 536--538). Every definition, display and label of those lines is retained unchanged. Omitted from this package and not redistributed: 1--531, 535--535, 539--569, 581--582, 647--1424. Nothing inside a retained excerpt is omitted or altered: every retained body is delivered exactly as the article's lines are written, so no omission receipt of any kind accompanies this package. Citations: the retained excerpts carry no citation command at all, so no bibliography entry of the article is transcribed and no reference is redistributed. Reference adaptation: none was needed and none was made; every reference inside the delivered unit resolves inside it, so no unresolved reference or citation is produced. Printed slips kept literal: no printed word, symbol, hypothesis or number of the counted statement or of its proof was altered. Layout and class: the article's own class is \texttt{svjour3} with packages including amsfonts, amsmath, graphicx, geometry, csquotes, amssymb, algorithm, algorithmicx, algpseudocode, color, hyperref, mathrsfs; the wrapper is the lane's pinned \texttt{amsart} editorial wrapper at 10pt on A4, which restates the theorem-like environments the retained text needs and adds the article's own macro definitions that the retained text needs (\texttt{\textbackslash P}, \texttt{\textbackslash de}), with extra wrapper packages (bm, bbm, dsfont, xspace, cancel, mathtools). The delivered numbering is the unit's own. Assets: no retained span contains a figure environment, a graphics-inclusion command, a PDF-page inclusion, a table or a TikZ picture; no image file is copied into this package, the package declares no asset dependency and no image file is redistributed with it. Licence: version arXiv:2506.03025v1 of this article is distributed under the Creative Commons Attribution 4.0 International licence, as recorded in the arXiv licence metadata for this exact version and shown on the abstract page https://arxiv.org/abs/2506.03025v1. A full rights notice is delivered at licenses/arxiv-2506.03025-CC-BY-4.0.txt. Characters: every retained excerpt is pure ASCII, so the delivered engine, XeLaTeX with its default Latin Modern text font, reports no missing character.
\begin{equation}\label{a1eq}
\hat{\boldsymbol{a}}_1 = R_{1}^{-1} \left(Q^T\boldsymbol{d} - R_{2} \hat{\boldsymbol{a}}_2\right).    
\end{equation}

\begin{equation}\label{a2eq}
	\hat{\boldsymbol{a}}_2=\left(A^{T} A\right)^{-1}A^{T}\boldsymbol{b}_1,
\end{equation}

\begin{proposition} \label{Thm:EstimateOfNorm}
Let $ \left\{ p_1, \ldots, p_D \right\} $ be a basis for $ \P_d\left(\mathbb{R}^2\right) $ with $\underset{i=1,\dots,D}{\max} \left\lVert p_i \right\rVert_{\infty} \leq 1$. Then,
\begin{equation}
\label{eq:BoundThm}
    \left\lVert \hat{\Pi}_d^{\mathrm{Pad}} \right\rVert_{\mathrm{op}} \leq \zeta_d + \eta_d,
\end{equation}
where
\begin{equation*}
\zeta_d \doteq \left\lVert R_{1}^{-1} \right\rVert_{1} \left(M \left\lVert Q^{T} \right\rVert_{1} + \left\lVert R_{2}\right\rVert_{1}\eta_d\right), \quad     \eta_d \doteq \left\lVert(A^T A)^{-1} A^T \right\rVert_1\left(D + M \left\lVert  W_1 R_{1}^{-1}Q^T\right\rVert_1  \right).
\end{equation*}
\end{proposition}

\begin{proof}
%Let $f\in C(\Omega)$ such that
%\begin{equation*}
%        \left\lVert f\right\rVert_{\infty}=\max_{\boldsymbol{x}\in \Omega}\left\lvert f(\boldsymbol{x})\right\rvert.
%\end{equation*}
%	We observe that
For any continuous function $f\in C(\Omega)$, we have
	\begin{equation} \label{eq:comoda}
%		\left\lVert \hat{\Pi}_d^{\mathrm{Pad}} \right\rVert_{\mathrm{op}} = \sup_{\Vert f \Vert_\infty = 1} 
\left\lVert \hat{\Pi}_d^{\mathrm{Pad}} [f] \right\rVert_\infty = \left\lVert \sum_{i=1}^D \hat{{a}}_i (f) p_i \right\rVert_\infty \leq \sum_{i=1}^D \left\lvert \hat{{a}}_i (f) \right\rvert = \left\Vert \hat{\boldsymbol{a}} (f) \right\Vert_1 = \left\Vert \hat{\boldsymbol{a}}_1 (f) \right\Vert_1 + \left\Vert \hat{\boldsymbol{a}}_2 (f) \right\Vert_1.
		\end{equation}
%Let $f\in C\left(\Omega\right)$ be such that   
%\begin{equation} \label{hyp}
%\left\lVert f \right\rVert_{\infty}\le 1.
%\end{equation}
% Let $\boldsymbol{x}\in\Omega$ and by using the triangular inequality, we obtain
%\begin{equation}
%\label{eq:ContinuityCMCLSO}
%    \left\lvert \hat{P}^{\mathrm{Pad}}_{r}[f](\boldsymbol{x}) \right\rvert=\left\lvert \sum_{i=1}^{\kappa(r)} \hat{a}_i(f)e_i(\boldsymbol{x}) \right\rvert\le \sum_{i=1}^{\kappa(r)} \left\lvert\hat{a}_i(f)\right\rvert \left\lvert e_i(\boldsymbol{x})\right\rvert.
%\end{equation}
%Then, by passing to the maximum to both sides, we have 
%\begin{equation} 
%\label{eq:BoundP(f)}
%   \left \lVert \hat{P}^{\mathrm{Pad}}_{r}[f] \right\rVert_{\infty} \le K\sum_{i=1}^{\kappa(r)} \left\lvert\hat{a}_i(f)\right\rvert=K\left\lVert\hat{\boldsymbol{a}}(f)\right\rVert_1,
%\end{equation}
%where $K$ is defined in~\eqref{eq:ConstantC}. Now, it remains to bound the discrete $L_1$-norm of   $\hat{\boldsymbol{a}}=\hat{\boldsymbol{a}}(f)$. We compute this vector by using the direct elimination method. Then, by~\eqref{vectdec}, we obtain
%\begin{eqnarray*}
%\left\lVert
%\hat{\boldsymbol{a}} 
%\right\rVert_1
%&=& 
%\left\lVert
%\begin{bmatrix}
%\hat{\boldsymbol{a}}_1 \\
%\hat{\boldsymbol{a}}_2
%\end{bmatrix}
%\right\rVert_1
%= \lVert \hat{\boldsymbol{a}}_1 \rVert_1 + \lVert \hat{\boldsymbol{a}}_2 \rVert_1.
%\end{eqnarray*}
Using the relations~\eqref{a1eq} and~\eqref{a2eq}, and applying the triangular inequality, we have
\begin{equation*}
    \left\lVert \hat{\boldsymbol{a}}_1 (f) \right\rVert_1\le \left\lVert R_{1}^{-1}\right\rVert_{1} \left(\left\lVert Q^T \right\rVert_1 \left\lVert \boldsymbol{d} (f) \right\rVert_1 + \left\lVert R_{2}\right\rVert_1 \left\lVert \hat{\boldsymbol{a}}_2 (f)\right\rVert_1\right) 
\end{equation*}
and
\begin{equation*}\label{a2}
   \left\lVert \hat{\boldsymbol{a}}_2 (f)\right\rVert_1\le\left\lVert\left(A^{T}A\right)^{-1}A^{T}\right\rVert_1 \left\lVert \boldsymbol{b}_1 (f) \right\rVert_1.
\end{equation*}
It remains to bound $\left\lVert \boldsymbol{b} (f) \right\rVert_1$ and $\left\lVert \boldsymbol{d} (f) \right\rVert_1$. By the mean value theorem, %Since $ \mathcal{T}_N$ is a partition of $ \Omega $, 
one has
\begin{equation*}\label{aus1}
    \left\lVert \boldsymbol{b} (f) \right\rVert_1=\sum_{i=1}^D \left\lvert \mu_i(f)\right\rvert\le\sum_{i=1} ^D \frac{1}{|t_i|}\int_{t_i}\left\lvert f(x,y)\right\rvert \de x \de y \le \left\Vert f \right\Vert_\infty \sum_{i=1} ^D 1 = D \left\Vert f \right\Vert_\infty
\end{equation*}
and 
\begin{equation*}
\left\lVert \boldsymbol{d} (f) \right\rVert_1=\sum_{i=1}^{M} \left\lvert \mu_i(f)\right\rvert\le \sum_{i=1}^M \frac{1}{|t_i|} \int_{t_i}\left\lvert f(x,y)\right\rvert \de x \de y \le \left\Vert f \right\Vert_\infty \sum_{i=1} ^{M} 1 = M \left\Vert f \right\Vert_\infty.
\end{equation*}
%Finally, by substituting~\eqref{B1} into~\eqref{a2}, applying the triangle inequality, and using~\eqref{aus1}, we obtain from~\eqref{eq:BoundP(f)}
%that
Since
\begin{equation*}
    \left\lVert \hat{\Pi}_d^{\mathrm{Pad}} \right\rVert_{\mathrm{op}} = \sup_{\left\Vert f \right\Vert_\infty \leq 1} \left\lVert \hat{\Pi}_d^{\mathrm{Pad}} [f] \right\rVert_\infty, 
\end{equation*}
plugging the above expressions into the bounds for $ \left\lVert \hat{\boldsymbol{a}}_1 (f) \right\rVert_1 $ and $ \left\lVert \hat{\boldsymbol{a}}_2 (f) \right\rVert_1 $, and taking the supremum over all $f \in C(\Omega)$ with $ \left\Vert f \right\Vert_\infty \leq 1 $, we get $ \left\lVert \hat{\boldsymbol{a}}_1 \right\rVert_1 \leq \zeta_d $ and $ \left\lVert \hat{\boldsymbol{a}}_2\right\rVert_1 \leq \eta_d $. Finally, substituting this into \eqref{eq:comoda}, we get the desired bound in \eqref{eq:BoundThm}.
\end{proof}

\end{document}
